% !PRINTED 588
\weberpage{588}
%<<CONT
und so bestimmt man eine Reihe ganzer Zahlen $\alpha, \alpha_1, \alpha_2, \alpha_3 \dots$
mit stets abnehmender Norm, die sich fortsetzen lässt, so lange
die Null nicht darunter vorkommt. Da diese Normen ganze rationale
positive Zahlen sind, die immer abnehmen, so muss nach
einer endlichen Zahl von Schritten die Null auftreten und man
erhält also ein Gleichungssystem

\weq{1}{\begin{aligned}
\alpha &= \mu\alpha_1 + \alpha_2,\\
\alpha_1 &= \mu_1\alpha_2 + \alpha_3,\\
&\cdots\cdots\cdots\cdots\\
\alpha_{h-2} &= \mu_{h-2}\alpha_{h-1} + \alpha_h,\\
\alpha_{h-1} &= \mu_{h-1}\alpha_h.
\end{aligned}}

Daraus schliesst man, dass die ganze Reihe der Zahlen
$\alpha, \alpha_1, \alpha_2 \dots$, also insbesondere auch $\alpha$ und $\alpha_1$ durch $\alpha_h$ theilbar
sind, und umgekehrt, dass jeder gemeinsamer Theiler von $\alpha$
und $\alpha_1$ Theiler von allen folgenden $\alpha$, also auch von $\alpha_h$ ist.
Jede andere Zahl, die diese beiden Eigenschaften hat, muss
mit $\alpha_h$ associirt sein, und man nennt also $\alpha_h$ (und jede mit $\alpha_h$
associirte Zahl) den \gesperrt{grössten gemeinschaftlichen Theiler} von
$\alpha$ und $\alpha_1$.

Wenn wir aus dem Algorithmus (1) die zwischenliegenden $\alpha$
eliminiren, indem wir $\alpha_2$ aus der ersten in die zweite, dann $\alpha_3$
aus der zweiten in die dritte Gleichung substituiren, so ergiebt
sich ein Resultat von der Form $\alpha\varkappa + \alpha_1\lambda = \alpha_h$, und wir können
also, wenn wir statt $\alpha, \alpha_1, \alpha_h$ setzen $\alpha, \beta, \delta$, den Satz aussprechen:

4. \gesperrt{Wenn $\alpha, \beta$ irgend zwei ganze Zahlen in $R(i)$
sind, und $\delta$ ihr grösster gemeinschaftlicher
Theiler, so kann man die ganzen Zahlen $\varkappa, \lambda$
so bestimmen, dass}

\weq{2}{\alpha\varkappa + \beta\lambda = \delta}
\gesperrt{wird.}

Ist der grösste gemeinschaftliche Theiler $\delta$ zweier Zahlen
$\alpha, \beta$ eine Einheit, so heissen $\alpha$ und $\beta$ relativ prim und wir
können in diesem Falle der Gleichung

\weq{3}{\alpha\varkappa + \beta\lambda = 1}
durch ganzzahlige $\varkappa, \lambda$ genügen.

Ist $\beta$ eine Primzahl, so ist entweder $\alpha$ durch $\beta$ theilbar,
oder $\alpha$ und $\beta$ sind relativ prim. Im letzteren Fall besteht die
Gleichung (3). Ist dann $\gamma$ eine andere Zahl in $R(i)$, so folgt aus (3)
\[
\alpha\gamma\varkappa + \beta\gamma\lambda = \gamma.
\]
